\chapter{Conclusion}
\label{ch:conclusion}

This thesis set out to build a difficulty estimator for classical guitar that a teacher could not only use but also interrogate. To this end, it adapted the additive RubricNet architecture from piano to classical guitar, guided by four research questions concerning descriptor design, predictive performance under an interpretability constraint, the multi-dimensional structure of difficulty, and the pedagogical validity of the model's explanations. The work makes four contributions:

\begin{enumerate}
    \item \textbf{A labeled symbolic dataset for classical guitar difficulty.} The dataset comprises 716 unique pieces pairing symbolic scores from three heterogeneous sources with professional 1--20 difficulty grades, assembled through fuzzy matching, manual auditing, and deduplication. Since the largest source (ClassClef PDF scores) carries no rhythm, a Needleman--Wunsch alignment pipeline transplants timing from matched MIDI performances onto the PDF-derived pitch sequences, leaving 89.2\% of the dataset with complete note durations.
    \item \textbf{A guitar-specific descriptor set, developed and audited across three generations.} Starting from 12 straightforward descriptors, correlation audits exposed uninformative features and extraction defects; the final set of 32 descriptors spans left-hand technique, right-hand technique, and global score complexity, and includes rhythm-aware descriptors --- windowed bottleneck measures and demand-per-beat proxies --- that encode the interaction between physical difficulty and available execution time.
    \item \textbf{The first adaptation of the additive RubricNet architecture beyond piano.} The final model reaches 31.0\% exact and 29.1\% balanced accuracy on the 8-class task (67.6\% on a coarse 3-level scale) and outperforms a Random Forest on balanced accuracy, MAE, MSE, and Kendall's $\tau$, while every prediction remains decomposable into per-descriptor contributions. Two fuzzy rule-based classifiers --- a complete search over primitive rules and fuzzy pattern trees, adapted here from music genre classification to difficulty grading --- serve as a second interpretable reference point; the model outperforms both on every metric, isolating the value of its ordinal design.
    \item \textbf{Evidence that the interpretability is substantive.} Descriptor contributions increase monotonically with difficulty, the model's influence rankings agree with Random Forest importances and model-free rank correlations, and the learned signs match pedagogical expectations --- open-string content lowers predicted difficulty, shift frequency raises it.
\end{enumerate}

Returning to the research questions: the descriptors that best capture graded difficulty are piece length, fretboard coverage and position statistics, shift frequency, and --- once rhythm is available --- demand-per-beat measures, with open-string content as the strongest negative indicator. A rubric-based model over these descriptors predicts difficulty with useful accuracy, matching a black-box baseline on ranking quality and surpassing it on error distance; equally interpretable fuzzy rule-based classifiers without ordinal awareness trail it by a wide margin on the same descriptors, showing that the combination of interpretability and ordinal design, not interpretability alone, is what keeps the cost of transparency near zero. Combining the three technical dimensions outperforms every single dimension alone, with the left hand strongest in isolation and each dimension contributing to the whole. Finally, the model's explanations align with pedagogical intuition, as demonstrated by the monotonicity and triangulation analyses.

A finding that runs through the entire project deserves restating: every substantial performance gain came from better data and better descriptors --- repairing a corrupted extraction path, recovering rhythm, encoding demand-versus-time interactions by hand --- while the model itself never changed. For interpretable-by-construction architectures, the descriptors are the model.

\section{Future Work}
\label{sec:future_work}
There are several avenues for future work.

One direction is to improve the treatment of pieces without rhythm information. The 10.8\% of pieces without rhythm currently receive training-fold median values for the eight rhythm-aware descriptors --- a safe but coarse fallback. A model trained on the rhythm-complete majority to predict rhythm descriptors from pitch-and-fret content could replace the medians with piece-specific estimates. A related and cheaper experiment worth conducting first is a re-evaluation on the rhythm-complete subset alone, to quantify exactly how much of V3's gain the imputed pieces dilute.

A second direction concerns the alignment pipeline's single-voice output, which discards notated voice structure. Splitting scores into melody and bass streams before alignment --- for example with the voice-separation utilities of the music21 toolkit~\cite{cuthbert2010music21} --- would preserve sustained bass durations and enable contrapuntal-complexity descriptors, an aspect of classical guitar difficulty the current set does not address.

Third, the descriptors currently trust the string-and-fret choices baked into each transcription. A guitar fingering solver that computes an ergonomic left-hand path from pitch content alone would let the descriptors measure a piece's \emph{best-case} physical difficulty rather than one transcriber's choices, and would additionally enable the grading of scores distributed as plain staff notation without tablature.

Fourth, the interaction-descriptor experiments showed that although guitar difficulty depends on multi-factor interactions, generic feature products only introduce collinearity into an additive model. A promising direction is decorrelated pairwise interactions, in the spirit of GA$^2$M~\cite{caruana2013intelligible}: rather than appending descriptor products directly, each candidate interaction could be residualized against its parent features, and the residuals --- statistically independent of the existing inputs --- fed to an otherwise-additive model. This would expand the interpretable model's capacity for interaction effects while preserving the rubric-like decomposition of its predictions.

Fifth, the fuzzy rule-based baselines invite an ordinal-aware variant. Their deficit was traced to the nominal one-vs-all decision rule rather than to the induced rules themselves, so decoding an expected grade from the per-class rule activations --- instead of taking their argmax --- could close much of the gap while retaining natural-language rules as the explanation format. The evolutionary rule-base induction of Heerde et al.~\cite{heerde2020comparing}, excluded here for its nondeterminism, likewise remains untested on difficulty grading.

Finally, all results in this thesis are relative to a single grading database. The natural next step is external validation: comparing the model's rubrics against independent grading systems, and placing the per-piece descriptor score sheets in front of practicing teachers to test whether the explanations change or support their judgments. The additive architecture makes such a study directly possible --- each prediction already constitutes an itemized rubric --- but whether teachers find the itemization convincing is an empirical question that remains open.
